\chapter{Aritmética dos Inteiros}\label{ch:b1:arith}

A aritmética --- o estudo da \omterm{def:b1:arith:divides}{divisibilidade} em $\Z$ --- foi iniciada no
volume do ensino médio. Este capítulo a reconstrói inteiramente a partir
da divisão euclidiana, com demonstrações completas: \omterm{thm:b1:arith:gcd}{máximo divisor comum}
e \omterm{met:b1:arith:euclid}{algoritmo de Euclides}, identidade de Bézout e lema de Gauss,
\omterm{thm:b1:arith:fta}{fatoração em primos} e o cálculo das \omterm{def:b1:arith:congruence}{congruências} até o pequeno teorema
de Fermat. Além do seu próprio encanto, este material é o modelo que o
\cref{ch:b1:poly} imita para os polinômios.

\section{Divisibilidade e divisão euclidiana}

\begin{definition}[Divisibilidade]\label{def:b1:arith:divides}
Para $a, b \in \Z$, $b$ \emph{divide} $a$ (escreve-se $b \mid
a$)\index{divisibilidade} quando $a = bq$ para algum $q \in \Z$.
Consequências básicas: se $b \mid a$ e $b \mid a'$, então $b \mid (ua + va')$
para todos $u, v \in \Z$; se $b \mid a$ e $a \neq 0$, então $\abs b \leq
\abs a$; e $a \mid b$ junto com $b \mid a$ forçam $b = \pm a$.
\end{definition}

\begin{theorem}[Divisão euclidiana]\label{thm:b1:arith:division}
Para todos $a \in \Z$ e $b \in \N^*$, existe exatamente um par $(q, r)
\in \Z \times \N$ com
\[
a = bq + r, \qquad 0 \leq r < b .
\]
\end{theorem}

\begin{proof}
\emph{Existência.} O \omterm{def:b1:logic:sets}{conjunto} $A = \{a - bk : k \in \Z\} \cap \N$ é um
subconjunto não vazio de $\N$ (tome $k = -\abs a$: $a + b\abs a \geq a +
\abs a \geq 0$). Seja $r = a - bq$ o seu menor elemento. Se $r \geq b$,
então $r - b = a - b(q+1)$ seria um elemento menor de $A$:
contradição. Logo, $0 \leq r < b$.

\emph{Unicidade.} Se $bq + r = bq' + r'$ com $0 \leq r, r' < b$,
então $b(q - q') = r' - r$ e $\abs{r' - r} < b$: o múltiplo de $b$
no lado esquerdo tem de ser $0$, de modo que $q = q'$ e $r = r'$.
\end{proof}

\begin{example}[Numeração posicional por divisões sucessivas]\label{ex:b1:arith:base7}
Escreva $2026$ na base $7$. Divida repetidamente por $7$, guardando os
restos:
\[
2026 = 7 \times 289 + 3, \quad
289 = 7 \times 41 + 2, \quad
41 = 7 \times 5 + 6, \quad
5 = 7 \times 0 + 5 .
\]
Lendo os restos do último ao primeiro: $2026 =
(5\,6\,2\,3)_7$. Verificação: $5 \times 343 + 6 \times 49 + 2 \times 7
+ 3 = 1715 + 294 + 14 + 3 = 2026$. A unicidade da divisão
euclidiana é exatamente o que torna cada algarismo \emph{forçado}: em
cada passo, o resto é o único inteiro de $\intint06$ congruente
ao valor corrente módulo $7$, de modo que a escrita na base $7$ é única
--- fato usado silenciosamente sempre que o problema de fim de semana
manipula ``os algarismos de $n$ na base $p$''.
\end{example}

\section{Máximo divisor comum}

\begin{theorem}[Subgrupos de $\Z$; existência do mdc]\label{thm:b1:arith:gcd}
\begin{enumerate}
  \item Todo subgrupo de $(\Z, +)$ é da forma $n\Z = \{nk : k
        \in \Z\}$ para um único $n \in \N$.
  \item Para $a, b \in \Z$ não simultaneamente nulos, o \omterm{def:b1:logic:sets}{conjunto} $a\Z + b\Z = \{au +
        bv : u, v \in \Z\}$ é um subgrupo de $\Z$ e, portanto, é igual a
        $d\,\Z$ para um único $d \in \N^*$. Esse $d$ é o
        \emph{máximo divisor comum}\index{máximo divisor comum}
        $\gcd(a, b)$: ele \omterm{def:b1:arith:divides}{divide} $a$ e $b$, e todo divisor comum de $a$ e
        $b$ \omterm{def:b1:arith:divides}{divide} $d$.
\end{enumerate}
\end{theorem}

\begin{proof}
(1) Seja $H \subseteq \Z$ um subgrupo (não vazio, estável por
subtração; a definição formal está no \cref{ch:b1:structures},
e apenas estas duas propriedades são usadas). Se $H = \{0\}$, tome $n = 0$.
Caso contrário, $H$ contém um elemento não nulo e o seu oposto, logo um
menor elemento estritamente positivo $n$. Então $n\Z \subseteq H$. Para $x
\in H$, escreva $x = nq + r$ com $0 \leq r < n$
(\cref{thm:b1:arith:division}); $r = x - nq \in H$, e a minimalidade de
$n$ força $r = 0$: $x \in n\Z$. Unicidade: $n$ é o menor elemento
positivo de $n\Z$.

(2) $a\Z + b\Z$ contém $0$ e é estável por subtração, de modo que é
$d\Z$ com $d \geq 1$ (ele contém $a$ ou $b$, não nulos). Como $a, b
\in d\Z$, $d$ \omterm{def:b1:arith:divides}{divide} ambos. E, se $c$ \omterm{def:b1:arith:divides}{divide} $a$ e $b$, então $c$
\omterm{def:b1:arith:divides}{divide} todo $au + bv$ --- em particular $c \mid d$, pois $d \in a\Z
+ b\Z$. Esta é a propriedade anunciada (e ela implica $\abs c \leq
d$, de modo que $d$ merece o nome de \emph{máximo} divisor comum).
\end{proof}

\begin{corollary}[Identidade de Bézout]\label{cor:b1:arith:bezout}
Para $a, b$ não simultaneamente nulos, existem $u, v \in \Z$ com
\[
au + bv = \gcd(a, b) .
\]
Em particular ($\gcd(a,b) = 1$, o caso dos números \emph{primos entre
si}\index{primos entre si}): $a$ e $b$ são \omterm{cor:b1:arith:bezout}{primos entre si} se, e somente
se, $au + bv = 1$ tem solução.
\end{corollary}

\begin{proof}
$\gcd(a,b) = d \in d\Z = a\Z + b\Z$. Quanto à equivalência: se
$\gcd(a,b) = 1$, Bézout fornece a solução; reciprocamente, $au + bv =
1$ força todo divisor comum de $a, b$ a dividir $1$.
\end{proof}

\begin{method}[Algoritmo de Euclides, estendido]\label{met:b1:arith:euclid}
Para calcular $\gcd(a, b)$ ($a > b > 0$): divida $a = bq + r$; então
$\gcd(a, b) = \gcd(b, r)$ (os divisores comuns de $(a,b)$ e os de $(b,r)$
coincidem, pois $r = a - bq$); itere até que o resto seja $0$; o
último resto não nulo é o mdc\index{algoritmo de Euclides}. Percorrendo
as divisões de trás para diante (ou mantendo os coeficientes na
descida), obtém-se um par de Bézout $(u, v)$.
\end{method}

\begin{example}\label{ex:b1:arith:euclid}
$\gcd(120, 23)$: $120 = 5 \times 23 + 5$; $23 = 4 \times 5 + 3$; $5 =
1\times 3 + 2$; $3 = 1 \times 2 + 1$; $2 = 2 \times 1 + 0$. Logo,
$\gcd = 1$. De trás para diante:
\begin{align*}
1 &= 3 - 2 = 3 - (5 - 3) = 2\times 3 - 5 = 2(23 - 4\times 5) - 5 \\
&= 2 \times 23 - 9 \times 5 = 2\times 23 - 9(120 - 5\times 23)
= 47 \times 23 - 9 \times 120 .
\end{align*}
Verificação: $47 \times 23 = 1081$, $9 \times 120 = 1080$.
\end{example}

\begin{theorem}[Lema de Gauss e consequências]\label{thm:b1:arith:gauss}
Sejam $a, b, c \in \Z$.
\begin{enumerate}
  \item (Lema de Gauss\index{lema de Gauss}) Se $a \mid bc$ e
        $\gcd(a, b) = 1$, então $a \mid c$.
  \item Se $a \mid c$, $b \mid c$ e $\gcd(a,b) = 1$, então $ab \mid
        c$.
  \item Se $\gcd(a, b) = \gcd(a, c) = 1$, então $\gcd(a, bc) = 1$.
\end{enumerate}
\end{theorem}

\begin{proof}
(1) Bézout: $au + bv = 1$. Multiplique por $c$: $acu + bcv = c$. Os dois
termos são divisíveis por $a$ (o segundo porque $a \mid bc$), logo $a
\mid c$.

(2) Escreva $c = aq$; de $b \mid aq$ e $\gcd(a, b) = 1$, o ponto (1)
dá $b \mid q$, de modo que $ab \mid aq = c$.

(3) $au + bv = 1$ e $au' + cv' = 1$. Multiplique as duas
relações:
\[
1 = (au + bv)(au' + cv')
= a\,\bigl(auu' + ucv' + u'bv\bigr) + bc\,(vv') ,
\]
uma relação de Bézout entre $a$ e $bc$: pelo
\cref{cor:b1:arith:bezout}, $\gcd(a, bc) = 1$.
\end{proof}

\begin{example}[Resolvendo uma equação diofantina linear]\label{ex:b1:arith:diophantine}
Encontre todos os $(x, y) \in \Z^2$ com $6x + 10y = 4$. Primeiro, o
\emph{teste de existência}: $\gcd(6, 10) = 2$ \omterm{def:b1:arith:divides}{divide} $4$, de modo que há
soluções (se o mdc não dividisse o lado direito, o lado esquerdo
seria sempre um múltiplo dele e não haveria nenhuma). Divida
tudo: $3x + 5y = 2$. Uma solução particular está à vista: $(x_0,
y_0) = (-1, 1)$. Para a geral, subtraia: $3(x + 1) = -5(y -
1)$, de modo que $3 \mid 5(y-1)$ e o lema de Gauss ($\gcd(3,5) = 1$) dá
$3 \mid y - 1$: $y = 1 - 3k$, e então $x = -1 + 5k$. Reciprocamente, todo
par desses serve:
\[
(x, y) = (-1 + 5k,\ 1 - 3k), \qquad k \in \Z .
\]
O padrão é geral: uma solução particular mais os múltiplos inteiros
de $\bigl(\frac b{\gcd}, -\frac a{\gcd}\bigr)$ --- a mesma estrutura
``particular mais homogênea'' do \cref{ch:b1:diffeq}, com o lema de Gauss
desempenhando o papel da unicidade.
\end{example}

\begin{definition}[Mínimo múltiplo comum]\label{def:b1:arith:lcm}
$\operatorname{lcm}(a, b)$\index{mínimo múltiplo comum} é o
gerador em $\N$ do subgrupo $a\Z \cap b\Z$: é um múltiplo comum de $a$ e
$b$ que \omterm{def:b1:arith:divides}{divide} todo múltiplo comum e, para $a, b \in \N^*$,
\[
\gcd(a,b) \times \operatorname{lcm}(a,b) = ab
\qquad (\text{demonstração no} \cref{exo:b1:arith:5}).
\]
\end{definition}

\begin{example}[Problemas de alinhamento são problemas de mmc]\label{ex:b1:arith:gears}
Duas engrenagens acopladas têm $84$ e $36$ dentes. Após quantos dentes
de movimento comum elas voltam juntas à posição inicial?
A configuração se repete quando o número de dentes decorridos é um
múltiplo comum de $84$ e $36$; a primeira vez é em
\[
\operatorname{lcm}(84, 36) = \frac{84 \times 36}{\gcd(84, 36)}
= \frac{3024}{12} = 252
\]
dentes --- isto é, $3$ voltas da engrenagem grande e $7$ da
pequena ($252/84$ e $252/36$). Note o caminho prático:
\emph{calcule primeiro o mdc} (Euclides: $84 = 2\times36 + 12$, $36
= 3\times12$) e depois divida --- nunca construa o mmc listando
múltiplos. Toda questão de coincidência periódica (engrenagens,
alinhamentos planetários, dízimas que se reencontram) reduz-se a esse
único cálculo.
\end{example}

\section{Números primos}

\begin{definition}\label{def:b1:arith:prime}
Um inteiro $p \geq 2$ é \emph{primo}\index{número primo} quando os
seus únicos divisores positivos são $1$ e $p$. Para $p$ primo e $a \in \Z$:
ou $p \mid a$, ou $\gcd(p, a) = 1$. Consequentemente
(\cref{thm:b1:arith:gauss}), vale o \emph{lema de Euclides}: se $p \mid
ab$, então $p \mid a$ ou $p \mid b$.
\end{definition}

\begin{remark}[Testando a primalidade por divisões sucessivas]\label{rem:b1:arith:trialdivision}
Se $n = ab$ com $2 \leq a \leq b$, então $a^2 \leq ab = n$, de modo que $a
\leq \sqrt n$: um $n$ composto tem sempre um divisor \omterm{def:b1:arith:prime}{primo} $\leq
\sqrt n$. Assim, para testar se $n$ é \omterm{def:b1:arith:prime}{primo} basta tentar
os \omterm{def:b1:arith:prime}{primos} até $\sqrt n$. Para $n = 271$: $\sqrt{271} < 17$, e
$271$ não é divisível por nenhum de $2, 3, 5, 7, 11, 13$ (é ímpar,
a soma dos algarismos é $10$, não termina em $0$ nem em $5$, $271 = 7\cdot38 + 5 =
11\cdot24 + 7 = 13\cdot20 + 11$): \omterm{def:b1:arith:prime}{primo}, após seis divisões
em vez de duzentas. A barreira $\sqrt n$ é um limiar genuíno:
cruzá-la eficientemente para números de cem algarismos exige
os testes modernos de primalidade nascidos do
\cref{thm:b1:arith:fermat}.
\end{remark}

\begin{theorem}[Euclides]\label{thm:b1:arith:euclidprimes}
Existem infinitos \omterm{def:b1:arith:prime}{números primos}.
\end{theorem}

\begin{proof}
Todo inteiro $n \geq 2$ tem um divisor \omterm{def:b1:arith:prime}{primo}: o seu menor divisor
$\geq 2$ é \omterm{def:b1:arith:prime}{primo} (uma fatoração própria dele produziria um
divisor menor de $n$). Agora suponha que $p_1, \dots, p_k$ fossem todos os
\omterm{def:b1:arith:prime}{primos} e seja $N = p_1 p_2 \cdots p_k + 1 \geq 2$. Algum \omterm{def:b1:arith:prime}{primo} $p_i$
\omterm{def:b1:arith:divides}{divide} $N$; mas $p_i$ também \omterm{def:b1:arith:divides}{divide} $N - 1 = p_1\cdots p_k$, de modo que $p_i
\mid 1$ --- absurdo.
\end{proof}

\begin{theorem}[Teorema fundamental da aritmética]\label{thm:b1:arith:fta}
Todo inteiro $n \geq 2$ é um produto de \omterm{def:b1:arith:prime}{primos}, e a fatoração
\[
n = p_1^{\alpha_1} p_2^{\alpha_2} \cdots p_k^{\alpha_k}
\qquad (p_1 < p_2 < \dots < p_k \text{ primos},\ \alpha_i \in \N^*)
\]
é única.\index{fatoração em primos}
\end{theorem}

\begin{proof}
\emph{Existência} por indução forte (\cref{thm:b1:logic:induction}):
$n = 2$ é \omterm{def:b1:arith:prime}{primo}; para $n > 2$, ou $n$ é \omterm{def:b1:arith:prime}{primo}, ou $n = ab$ com
$2 \leq a, b < n$, e a hipótese de indução fatora $a$ e $b$.

\emph{Unicidade.} Suponha $p_1 \cdots p_r = q_1 \cdots q_s$
(\omterm{def:b1:arith:prime}{primos} listados com repetição, digamos $r \leq s$) e faça indução em
$r$. Se $r = 0$, o lado esquerdo é $1$, o que força $s = 0$ (um produto
não vazio de \omterm{def:b1:arith:prime}{primos} excede $1$). Para $r \geq 1$: o \omterm{def:b1:arith:prime}{primo} $p_1$
\omterm{def:b1:arith:divides}{divide} $q_1(q_2\cdots q_s)$, de modo que, pelo lema de Euclides, ou $p_1
\mid q_1$, ou $p_1 \mid q_2\cdots q_s$; iterando, $p_1$ \omterm{def:b1:arith:divides}{divide}
algum $q_j$. Mas $q_j$ é \omterm{def:b1:arith:prime}{primo} e $p_1 \geq 2$: necessariamente $p_1
= q_j$. Cancele esse fator comum (legítimo: $\Z$ é um
domínio de integridade) para obter
\[
p_2 \cdots p_r = q_1 \cdots \widehat{q_j} \cdots q_s
\]
(com o chapéu marcando a omissão), uma igualdade de produtos mais
curtos; a hipótese de indução diz que as duas listas $p_2, \dots, p_r$ e
$q_1, \dots, \widehat{q_j}, \dots, q_s$ coincidem a menos da ordem
e, portanto, as originais também coincidiam. A forma com expoentes
agrupa os \omterm{def:b1:arith:prime}{primos} iguais.
\end{proof}

\begin{proposition}[Valorações]\label{prop:b1:arith:valuation}
Para $p$ \omterm{def:b1:arith:prime}{primo} e $n \in \N^*$, escreva $v_p(n)$\index{valoracao
p-adica@valoração $p$-ádica} para o expoente de $p$ na
fatoração de $n$ (com $v_p(n) = 0$ se $p \nmid n$). Então
\[
v_p(mn) = v_p(m) + v_p(n),
\qquad
m \mid n \iff \forall p,\ v_p(m) \leq v_p(n),
\]
\[
v_p\bigl(\gcd(m,n)\bigr) = \min\bigl(v_p(m), v_p(n)\bigr),
\qquad
v_p\bigl(\operatorname{lcm}(m,n)\bigr) = \max\bigl(v_p(m),
v_p(n)\bigr).
\]
\end{proposition}

\begin{proof}
A primeira identidade vale porque as fatorações se multiplicam e a
fatoração de $mn$ é única. Se $m \mid n$, escreva $n = mq$ e
aplique-a. Reciprocamente, se todos os $v_p(m) \leq v_p(n)$, o inteiro $q =
\prod_p p^{\,v_p(n) - v_p(m)}$ satisfaz $mq = n$. A fórmula do mdc:
o inteiro $d = \prod p^{\min}$ \omterm{def:b1:arith:divides}{divide} os dois pelo critério, e
todo divisor comum $c$ tem $v_p(c) \leq \min$ para todo $p$, de modo que $c
\mid d$; mesmo raciocínio para o mmc com o $\max$.
\end{proof}

\begin{example}[Quadrados e cubos por meio de valorações]\label{ex:b1:arith:squarecube}
Um inteiro $n \geq 1$ é um quadrado perfeito se, e somente se, todo
$v_p(n)$ é par (se $n = m^2$, então $v_p(n) = 2v_p(m)$;
reciprocamente, divida ao meio cada expoente). Do mesmo modo para os
cubos, com múltiplos de $3$. Assim, $21168 = 2^4 \times 3^3 \times 7^2$ não é
um quadrado ($v_3 = 3$ é ímpar) nem um cubo ($v_2 = 4$); o
menor inteiro positivo $m$ tal que $21168\,m$ \emph{seja} um
cubo é encontrado completando cada expoente até o próximo múltiplo de
$3$:
\[
m = 2^{6-4} \times 3^{3-3} \times 7^{3-2} = 2^2 \times 7 = 28,
\qquad
21168 \times 28 = 2^6\,3^3\,7^3 = (2^2 \times 3 \times 7)^3
= 84^3 .
\]
A ideia: questões multiplicativas (quadrados, cubos, divisores,
mdc, mmc) tornam-se questões \emph{coordenada a coordenada} sobre os
vetores de expoentes $(v_2, v_3, v_5, \dots)$ --- a fatoração única é a
afirmação de que essas coordenadas existem e estão bem definidas.
\end{example}

\section{Congruências}

\begin{definition}\label{def:b1:arith:congruence}
Para $n \in \N^*$: $a \equiv b \pmod n$\index{congruência} quando $n \mid
a - b$. Esta é uma \omterm{def:b1:logic:equiv}{relação de equivalência} compatível com a adição e
a multiplicação: se $a \equiv b$ e $a' \equiv b'$ (mód.\ $n$), então
$a + a' \equiv b + b'$, $aa' \equiv bb'$ e $a^k \equiv b^k$ para
$k \in \N$.
\end{definition}

\begin{example}[A prova dos noves]\label{ex:b1:arith:castingnines}
A compatibilidade com $+$ e $\times$ é um recurso de verificação tão
antigo quanto o comércio. Como $10 \equiv 1 \pmod 9$, todo inteiro é
congruente módulo $9$ à soma dos seus algarismos (demonstrado no
\cref{exo:b1:arith:2}). Para verificar a afirmação $1234 \times 567 =
699\,678$: as somas dos algarismos dão $1234 \equiv 1$ e $567 \equiv 18
\equiv 0 \pmod 9$, de modo que o produto deve ser $\equiv 1 \times 0 =
0$; e, de fato, $6 + 9 + 9 + 6 + 7 + 8 = 45 \equiv 0$. A verificação
passa (e o produto está, de fato, correto). Se alguém tivesse relatado
$699\,478$, a soma dos algarismos $43 \equiv 7 \not\equiv 0$ o
condenaria instantaneamente. O teste é unilateral --- ele apanha um
erro a menos que o próprio erro seja um múltiplo de $9$ --- o que é
exatamente a lição dos pseudoprimos do
\cref{ex:b1:arith:pseudoprime} em miniatura: verificações por \omterm{def:b1:arith:congruence}{congruência}
refutam, não certificam.
\end{example}

\begin{proposition}[Invertibilidade módulo $n$]\label{prop:b1:arith:invmod}
$a$ é \emph{invertível módulo $n$} (isto é, $ab \equiv 1 \pmod n$ para algum
$b$) se, e somente se, $\gcd(a, n) = 1$. A inversa é então única módulo
$n$ e calculada pelo \omterm{met:b1:arith:euclid}{algoritmo de Euclides} estendido.
\end{proposition}

\begin{proof}
$ab \equiv 1 \pmod n$ significa $ab + nk = 1$ para algum $k$: uma relação
de Bézout, que existe se, e somente se, $\gcd(a,n) = 1$
(\cref{cor:b1:arith:bezout}). Unicidade: se $ab \equiv ab' \equiv 1$,
então $b \equiv b(ab') = (ab)b' \equiv b' \pmod n$.
\end{proof}

\begin{example}[Invertendo $7$ módulo $26$]\label{ex:b1:arith:inv7mod26}
Como $\gcd(7, 26) = 1$, a classe de $7$ é invertível módulo $26$.
Euclides estendido:
\[
26 = 3 \times 7 + 5, \qquad
7 = 1 \times 5 + 2, \qquad
5 = 2 \times 2 + 1 ,
\]
e depois, de trás para diante:
\[
1 = 5 - 2 \times 2 = 5 - 2(7 - 5) = 3 \times 5 - 2 \times 7
= 3(26 - 3 \times 7) - 2 \times 7 = 3 \times 26 - 11 \times 7 .
\]
Portanto, $7 \times (-11) \equiv 1 \pmod{26}$, isto é, $7^{-1} \equiv
-11 \equiv 15 \pmod{26}$; verificação: $7 \times 15 = 105 = 4 \times 26
+ 1$. Com a inversa em mãos, qualquer \omterm{def:b1:arith:congruence}{congruência} $7x \equiv c
\pmod{26}$ resolve-se com uma multiplicação: $x \equiv 15c$. Essa
inversão mecânica é o cavalo de batalha da aritmética modular ---
e dos protocolos de chave pública mencionados na
\cref{rem:b1:arith:whereused}, em que os \omterm{def:b1:complex:field}{módulos} têm centenas de
algarismos, mas o algoritmo é exatamente este.
\end{example}

\begin{example}[Quando o coeficiente não é invertível]\label{ex:b1:arith:noninvertible}
Resolva $12x \equiv 8 \pmod{20}$. Aqui $\gcd(12, 20) = 4$, de modo que $12$
não é invertível módulo $20$ --- mas a equação ainda é
tratável. A \omterm{def:b1:arith:congruence}{congruência} diz que $20 \mid 12x - 8$; dividindo a
relação inteira por $4$ (divisor dos três ingredientes), ela é
equivalente a $5 \mid 3x - 2$, isto é,
\[
3x \equiv 2 \pmod 5 .
\]
Agora, $\gcd(3, 5) = 1$ e $3^{-1} \equiv 2 \pmod 5$ ($3 \times 2 =
6 \equiv 1$), de modo que $x \equiv 4 \pmod 5$: as soluções são $x
\equiv 4, 9, 14, 19 \pmod{20}$ --- \emph{quatro} classes módulo $20$,
correspondendo ao mdc. (Se o lado direito não fosse divisível por $4$,
digamos $12x \equiv 6 \pmod{20}$, não haveria solução alguma:
o lado esquerdo é sempre $\equiv 0 \pmod 4$.) Forma geral: $ax
\equiv b \pmod n$ tem solução se, e somente se, $\gcd(a, n) \mid b$, e nesse caso
tem exatamente $\gcd(a, n)$ classes de soluções --- divida tudo
pelo mdc e inverta.
\end{example}

\begin{theorem}[Pequeno teorema de Fermat]\label{thm:b1:arith:fermat}
Seja $p$ \omterm{def:b1:arith:prime}{primo}. Para todo $a \in \Z$:
\[
a^p \equiv a \pmod p,
\]
e, se $p \nmid a$, então $a^{p-1} \equiv 1 \pmod
p$.\index{pequeno teorema de Fermat}
\end{theorem}

\begin{proof}
Primeiro, para $1 \leq k \leq p - 1$, o \omterm{def:b1:counting:objects}{coeficiente binomial} $\binom pk
= \frac{p!}{k!(p-k)!}$ é divisível por $p$: com efeito, $k!\,(p-k)!\,
\binom pk = p!$ e $p$ \omterm{def:b1:arith:divides}{divide} $p!$, mas é \omterm{def:b1:arith:prime}{primo} com $k!(p-k)!$
(todos os fatores são $< p$), de modo que o lema de Gauss dá $p \mid \binom pk$.

Agora demonstremos $a^p \equiv a$ para $a \in \N$ por indução. Verdadeiro para $a =
0$. Se $a^p \equiv a$, então, pelo teorema binomial,
\[
(a+1)^p = \sum_{k=0}^{p} \binom pk a^k
\equiv a^p + 1 \equiv a + 1 \pmod p,
\]
com todos os termos intermediários se anulando módulo $p$. Para $a < 0$,
aplique o resultado a $-a$ e separe $p = 2$ (em que $x \equiv -x$) do $p$ ímpar
(em que $(-a)^p = -a^p$). Por fim, se $p \nmid a$, multiplique $a^p \equiv a$ por
uma inversa de $a$ módulo $p$ (\cref{prop:b1:arith:invmod}).
\end{proof}

\begin{example}[A recíproca de Fermat falha: $341$]\label{ex:b1:arith:pseudoprime}
O pequeno teorema de Fermat dá um teste barato de
\emph{composicionalidade}: se $a^{n-1} \not\equiv 1 \pmod n$ para algum $a$ \omterm{def:b1:arith:prime}{primo} com $n$,
então $n$ não é \omterm{def:b1:arith:prime}{primo}. O teste poderia também certificar a primalidade?
Não: tome $n = 341 = 11 \times 31$, composto, e $a = 2$. Como
$2^{10} = 1024 = 3 \times 341 + 1$,
\[
2^{10} \equiv 1 \pmod{341}
\qquad\Longrightarrow\qquad
2^{340} = \bigl(2^{10}\bigr)^{34} \equiv 1 \pmod{341} :
\]
o composto $341$ passa no teste de Fermat na base $2$ (é o menor
\emph{pseudoprimo} desse tipo). A base $3$ o desmascara
($3^{340} \not\equiv 1$), e o teste prático de primalidade,
portanto, roda o teste em várias bases, com refinamentos ---
as versões industriais dessa ideia são o que certifica os grandes
\omterm{def:b1:arith:prime}{primos} da \cref{rem:b1:arith:whereused}. Moral: uma implicação
e a sua recíproca vivem vidas separadas
(\cref{rem:b1:logic:pitfalls}), mesmo para teoremas.
\end{example}

\begin{example}[Cálculos práticos com congruências]\label{ex:b1:arith:congruences}
Qual é o resto de $7^{2026}$ na divisão por $11$? Por Fermat, $7^{10}
\equiv 1 \pmod{11}$. Como $2026 = 10 \times 202 + 6$:
\[
7^{2026} \equiv 7^6 = (7^2)^3 = 49^3 \equiv 5^3 = 125 \equiv 4
\pmod{11}.
\]
O resto é $4$. A estratégia: reduza o expoente \omterm{def:b1:complex:field}{módulo} a
ordem fornecida por Fermat e depois reduza as potências intermediárias a
cada passo.
\end{example}

\begin{remark}[Armadilhas frequentes em aritmética]\label{rem:b1:arith:pitfalls}
\begin{enumerate}
  \item \emph{Dividir uma \omterm{def:b1:arith:congruence}{congruência}.} De $ac \equiv bc \pmod n$
        \emph{não} se pode concluir $a \equiv b$, a menos que $\gcd(c,
        n) = 1$: $6 \equiv 2 \pmod 4$, mas $3 \not\equiv 1 \pmod
        4$. A regra geral correta \omterm{def:b1:arith:divides}{divide} também o \omterm{def:b1:complex:field}{módulo}: $ac
        \equiv bc \pmod n \iff a \equiv b \pmod{n/\gcd(c,n)}$.
  \item \emph{Usar mal o lema de Euclides.} $a \mid bc$ implica $a
        \mid b$ ou $a \mid c$ apenas para $a$ \emph{\omterm{def:b1:arith:prime}{primo}} (ou
        \omterm{def:b1:arith:prime}{primo} com um dos fatores): $6 \mid 4 \times 9$, mas $6$
        não \omterm{def:b1:arith:divides}{divide} nenhum dos fatores.
  \item \emph{``\omterm{cor:b1:arith:bezout}{Primos entre si}'' é uma relação, não uma propriedade.}
        ``$8$ e $9$ são \omterm{cor:b1:arith:bezout}{primos entre si}'' é verdade, embora nenhum
        deles seja \omterm{def:b1:arith:prime}{primo}; ``dois a dois \omterm{cor:b1:arith:bezout}{primos entre si}'' é mais forte
        que ``\omterm{cor:b1:arith:bezout}{primos entre si} no \omterm{def:b1:logic:sets}{conjunto}'' ($\gcd(6, 10, 15) = 1$, mas
        nenhum par é formado por \omterm{cor:b1:arith:bezout}{primos entre si}).
  \item \emph{Os expoentes não vivem módulo $n$.} Em $a^k \bmod n$,
        o expoente só pode ser reduzido \omterm{def:b1:complex:field}{módulo} a \emph{ordem}
        de $a$ (por exemplo, $p - 1$ quando Fermat se aplica), nunca
        módulo $n$: $2^{10} \bmod 11$ vale $1$, e não $2^{10 \bmod
        11} = 2^{10}$ --- a redução que funciona é a que o
        \cref{ex:b1:arith:congruences} executa.
\end{enumerate}
\end{remark}

\begin{remark}[Onde este capítulo é usado]\label{rem:b1:arith:whereused}
Este capítulo é tanto um modelo quanto uma caixa de ferramentas. Toda a
cadeia --- divisão euclidiana, mdc, Bézout, Gauss, fatoração única ---
é reencenada literalmente para os polinômios no \cref{ch:b1:poly},
em que o ``grau'' faz o papel do valor absoluto; comparar os
dois capítulos lado a lado é a melhor maneira de entender ambos. O
cálculo das \omterm{def:b1:arith:congruence}{congruências} torna-se o anel $\Z/n\Z$ no
\cref{ch:b1:structures}, cujos elementos invertíveis
(\cref{prop:b1:arith:invmod}) formam o primeiro exemplo não trivial de
um grupo de unidades. As valorações voltam no problema de fim de semana
abaixo (fórmula de Legendre) e alimentam as demonstrações de
irracionalidade do \cref{ch:b1:reals}. Além deste volume, a inversão de
Bézout módulo $n$ é o motor da criptografia de chave pública, e o
pequeno teorema de Fermat é o avô dos testes de primalidade que
certificam os grandes \omterm{def:b1:arith:prime}{primos} ali usados.
\end{remark}

\begin{remark}[Interlúdio: $\Z$ como modelo]\label{rem:b1:arith:interlude}
Afaste-se dos teoremas individuais e observe a
arquitetura do capítulo: uma ferramenta (a divisão euclidiana)
produziu uma classificação (os subgrupos $n\Z$), que produziu um
teorema de existência (mdc, Bézout), que produziu um cálculo de
\omterm{def:b1:arith:divides}{divisibilidade} (Gauss), que produziu a fatoração única --- cada
andar apoiado apenas no de baixo. O mesmo edifício será
erguido mais duas vezes neste volume, com térreos diferentes:
no \cref{ch:b1:poly}, em que a divisão pelo grau substitui a divisão
pelo tamanho e tudo acima se repete \emph{literalmente}; e, em
miniatura, dentro de cada $\Z/n\Z$ do \cref{ch:b1:structures},
em que questões de invertibilidade (a \cref{prop:b1:arith:invmod}
deste capítulo) se tornam enunciados estruturais sobre
anéis e corpos. Reconhecer um argumento como ``o argumento de $\Z$,
transplantado'' é o modo mais rápido de aprender esses capítulos --- e
o primeiro sabor do hábito central da álgebra: demonstrar teoremas
sobre \emph{axiomas}, e não sobre objetos.
\end{remark}

\begin{omfigure}
\begin{tikzpicture}[scale=0.42]
  \foreach \n/\k in {0/0, 1/0, 1/1, 2/0, 2/2, 3/0, 3/1, 3/2, 3/3,
      4/0, 4/4, 5/0, 5/1, 5/4, 5/5, 6/0, 6/2, 6/4, 6/6,
      7/0, 7/1, 7/2, 7/3, 7/4, 7/5, 7/6, 7/7}
    \fill[omDef] ({2*\k - \n}, {-\n}) circle (0.32);
  \foreach \n/\k in {2/1, 4/1, 4/2, 4/3, 5/2, 5/3, 6/1, 6/3, 6/5}
    \draw[gray] ({2*\k - \n}, {-\n}) circle (0.32);
\end{tikzpicture}

{\small Linhas $0$ a $7$ do triângulo de Pascal com as entradas
\emph{ímpares} preenchidas: a linha $n$ contém $2^{s_2(n)}$ delas, em que
$s_2(n)$ é o número de uns na escrita binária de $n$
(linhas $1, 2, 4$: duas entradas ímpares; linha $7 = (111)_2$: todas as oito).
O padrão autossemelhante --- cada ``triângulo de ímpares'' gera duas
cópias de si mesmo --- é o teorema de Kummer em forma de figura,
demonstrado no problema de fim de semana abaixo.}
\end{omfigure}

\section{Exercícios}

\begin{exercise}[$\star$]\label{exo:b1:arith:1}
Calcule $\gcd(1\,001, 777)$ pelo \omterm{met:b1:arith:euclid}{algoritmo de Euclides}, e um par de
Bézout para ele.
\end{exercise}

\begin{exercise}[$\star$]\label{exo:b1:arith:2}
Demonstre os critérios de \omterm{def:b1:arith:divides}{divisibilidade} na base $10$: um inteiro é
congruente módulo $9$ à soma dos seus algarismos, e módulo $11$ à soma
alternada dos seus algarismos. Quanto vale $123\,456\,789$ módulo $9$ e
módulo $11$?
\end{exercise}

\begin{exercise}[$\star$]\label{exo:b1:arith:3}
Resolva em $\Z$: $91x \equiv 1 \pmod{237}$ \emph{(Euclides estendido)}.
\end{exercise}

\begin{exercise}[$\star$]\label{exo:b1:arith:4}
Encontre todos os pares $(x, y) \in \Z^2$ com $17x + 39y = 1$; depois todos os
pares com $17 x + 39 y = 5$.
\end{exercise}

\begin{exercise}[$\star\star$]\label{exo:b1:arith:5}
Demonstre que, para $a, b \in \N^*$: $\gcd(a,b) \times
\operatorname{lcm}(a,b) = ab$. \emph{(Use as fórmulas de valoração da
\cref{prop:b1:arith:valuation} e $\min(\alpha,\beta) +
\max(\alpha,\beta) = \alpha + \beta$.)}
\end{exercise}

\begin{exercise}[$\star\star$]\label{exo:b1:arith:6}
Sejam $a = 2^{10} \times 3^4 \times 5^2$ e $b = 2^6 \times 3^7 \times
7$. Calcule $\gcd(a, b)$, $\operatorname{lcm}(a,b)$ e o número de
divisores positivos de $a$. \emph{(Demonstre a fórmula de contagem de
divisores $\prod_i (\alpha_i + 1)$.)}
\end{exercise}

\begin{exercise}[$\star\star$]\label{exo:b1:arith:7}
Demonstre que $\sqrt p$ é irracional para todo \omterm{def:b1:arith:prime}{primo} $p$, usando
valorações: compare $v_p$ dos dois lados de $p q^2 = r^2$.
\end{exercise}

\begin{exercise}[$\star\star$]\label{exo:b1:arith:8}
(Problema chinês do resto) Encontre todos os inteiros $x$ com
\[
x \equiv 2 \pmod 7, \qquad x \equiv 5 \pmod{11}.
\]
Demonstre, ao longo do caminho, que, para $m, n$ \omterm{cor:b1:arith:bezout}{primos entre si}, o par
de \omterm{def:b1:arith:congruence}{congruências} $x \equiv a \ (m)$, $x \equiv b\ (n)$ sempre tem solução, única
módulo $mn$.\index{teorema chinês do resto}
\end{exercise}

\begin{exercise}[$\star\star$]\label{exo:b1:arith:9}
Calcule $3^{1000}$ módulo $7$ e os dois últimos algarismos decimais de
$7^{100}$ \emph{(módulo $100 = 4 \times 25$: use o
\cref{exo:b1:arith:8})}.
\end{exercise}

\begin{exercise}[$\star\star\star$]\label{exo:b1:arith:10}
Para $m, n \in \N^*$, demonstre que $\gcd(2^m - 1,\, 2^n - 1) =
2^{\gcd(m,n)} - 1$. \emph{Sugestão: mostre primeiro que o resto de $2^m
- 1$ módulo $2^n - 1$ é $2^r - 1$, em que $r$ é o resto de $m$ módulo
$n$; depois siga o \omterm{met:b1:arith:euclid}{algoritmo de Euclides}.}
\end{exercise}

\begin{exercise}[$\star\star\star$]\label{exo:b1:arith:11}
(Teorema de Wilson) Seja $p$ um \omterm{def:b1:arith:prime}{primo}. Demonstre que
\[
(p-1)! \equiv -1 \pmod p ,
\]
emparelhando cada fator de $(p-1)!$ com a sua inversa módulo $p$ e
identificando os fatores emparelhados consigo mesmos (resolva antes
$x^2 \equiv 1 \pmod p$). Verifique a recíproca: se $n \geq 2$ não é \omterm{def:b1:arith:prime}{primo}, então $(n-1)!
\not\equiv -1 \pmod n$.
\end{exercise}

\begin{exercise}[$\star\star\star$]\label{exo:b1:arith:12}
(Números de Fermat) Para $n \in \N$, seja $F_n = 2^{2^n} + 1$.
\begin{enumerate}
  \item Demonstre que $F_0 F_1 \cdots F_{n-1} = F_n - 2$ para $n \geq
        1$ (indução).
  \item Deduza que os números de Fermat são dois a dois \omterm{cor:b1:arith:bezout}{primos entre si}.
  \item Deduza uma segunda demonstração, independente do
        \cref{thm:b1:arith:euclidprimes}, de que existem
        infinitos \omterm{def:b1:arith:prime}{primos}.
\end{enumerate}
\end{exercise}

\section{Problema: A fórmula de Legendre e os transportes de Kummer}

\begin{problem}\label{pb:b1:arith:1}
Quantos zeros terminam a escrita decimal de $1000!$ --- e, mais a
fundo, qual é a potência exata de um \omterm{def:b1:arith:prime}{primo} $p$ que \omterm{def:b1:arith:divides}{divide} $n!$ ou que
\omterm{def:b1:arith:divides}{divide} um \omterm{def:b1:counting:objects}{coeficiente binomial}? As respostas completas são duas joias
da aritmética elementar: a \emph{fórmula de Legendre}
\index{fórmula de Legendre} $v_p(n!) = \sum_{k\geq1}
\lfloor n/p^k \rfloor$, com o seu avatar digital $v_p(n!) = \frac{n
- s_p(n)}{p-1}$, e o \emph{teorema de Kummer}\index{teorema de
Kummer}: $v_p\binom{m+n}m$ conta os \emph{transportes} (os ``vai um'') na
soma de $m$ e $n$ na base $p$. Este problema demonstra os dois, confere
um contra o outro numericamente e colhe as consequências clássicas ---
zeros finais, a paridade do triângulo de Pascal e uma primeira
estimativa na direção do teorema dos \omterm{def:b1:arith:prime}{números primos}.
Ao longo do problema, $p$ é um \omterm{def:b1:arith:prime}{primo}, $\floor{x}$ é a parte inteira e
$s_p(n)$ denota a soma dos algarismos de $n$ escrito na base $p$.

\medskip
\noindent\textbf{Parte I --- Partes inteiras, valorações e a fórmula
de Legendre.}
\begin{enumerate}
  \item Aquecimento: calcule $10!$ e leia o seu número de zeros
        finais; calcule $v_2(10!)$ e $v_5(10!)$ diretamente a partir da
        fatoração de cada fator $1, 2, \dots, 10$.
  \item Demonstre que, para $x \in \R$ e $n \in \N^*$,
        $\bigl\lfloor \lfloor x \rfloor / n \bigr\rfloor =
        \lfloor x/n \rfloor$.
  \item Demonstre que $v_p(a + b) \geq \min\bigl(v_p(a),
        v_p(b)\bigr)$ para todos $a, b \in \N^*$, com igualdade
        sempre que $v_p(a) \neq v_p(b)$.
  \item Mostre que o número de múltiplos de $m$ em $\intint1n$ é
        $\lfloor n/m \rfloor$.
  \item Demonstre a \emph{fórmula de Legendre}: para todo $n \in \N^*$,
        \[
        v_p(n!) = \sum_{k=1}^{\infty}
        \Bigl\lfloor \frac{n}{p^k} \Bigr\rfloor
        \]
        (uma soma finita: os termos se anulam assim que $p^k > n$).
        \emph{Conte, para cada $k$, os fatores de $\intint1n$
        divisíveis por $p^k$: cada um contribui com exatamente uma
        unidade por nível que alcança.}
\end{enumerate}

\medskip
\noindent\textbf{Parte II --- A forma digital e os zeros
finais.}
\begin{enumerate}
  \item Calcule $v_5(1000!)$ e $v_2(1000!)$, e conclua: quantos
        zeros terminam $1000!$?
  \item Demonstre a forma digital da fórmula de Legendre: escrevendo
        $n =
        \sum_i a_i p^i$ na base $p$,
        \[
        v_p(n!) = \frac{n - s_p(n)}{p - 1} .
        \]
  \item Duas consequências para $p = 2$: mostre que $2^n$ nunca
        \omterm{def:b1:arith:divides}{divide} $n!$, e que $2^{n-1}$ \omterm{def:b1:arith:divides}{divide} $n!$ exatamente quando
        $n$ é uma potência de $2$.
  \item Estime o defeito: mostre que $\frac n{p-1} - \log_p(n) - 1 \leq
        v_p(n!) < \frac n{p-1}$, de modo que $\frac{v_p(n!)}{n} \to
        \frac1{p-1}$: a longo prazo, acumula-se uma proporção
        $\frac1{p-1}$ de um fator $p$ por unidade.
  \item Seja $Z(n) = v_5(n!)$ o número de zeros finais de
        $n!$. Mostre que $Z(n) - Z(n-1) = v_5(n)$, deduza que $Z$ pula
        inteiramente o valor $5$ (calcule $Z(24)$ e $Z(25)$) e
        demonstre que nenhum fatorial termina em exatamente cinco zeros.
\end{enumerate}

\medskip
\noindent\textbf{Parte III --- O teorema de Kummer.}
\begin{enumerate}
  \item Demonstre que $\lfloor x + y \rfloor - \lfloor x \rfloor -
        \lfloor y \rfloor \in \{0, 1\}$ para todos $x, y \in \R$, e
        deduza da fórmula de Legendre que
        \[
        v_p\binom{m+n}m
        = \sum_{k\geq1}\Bigl(
        \Bigl\lfloor\frac{m+n}{p^k}\Bigr\rfloor
        - \Bigl\lfloor\frac{m}{p^k}\Bigr\rfloor
        - \Bigl\lfloor\frac{n}{p^k}\Bigr\rfloor\Bigr),
        \]
        uma soma cujos termos valem cada um $0$ ou $1$.
  \item Demonstre o \emph{teorema de Kummer}: o $k$-ésimo termo dessa
        soma vale $1$ exatamente quando a soma de $m$ e $n$ na
        base $p$ produz um transporte para a posição $k$; portanto,
        $v_p\binom{m+n}m$ é o número total de transportes.
        \emph{(Escreva $m = p^km_1 + m_0$ e $n = p^kn_1 + n_0$
        com $0 \leq m_0, n_0 < p^k$ e examine $\lfloor (m_0 +
        n_0)/p^k \rfloor$.)}
  \item Deduza que, para $0 < j < p^k$:
        \[
        v_p\binom{p^k}{j} = k - v_p(j) ,
        \]
        contando os transportes na soma $j + (p^k - j)$.
        (Em particular, $p \mid \binom p j$ para $0 < j < p$: o
        passo-chave do \cref{thm:b1:arith:fermat}, recuperado.)
  \item Demonstre que $v_2\binom{2n}n = s_2(n)$. Deduza que o
        \omterm{def:b1:counting:objects}{coeficiente binomial} central é sempre par e que
        $\binom{2n}n \equiv 2 \pmod 4$ exatamente quando $n$ é uma potência
        de $2$.
  \item Mostre, usando a identidade de Vandermonde
        (\cref{exo:b1:counting:7}) e a questão 13, que
        $\binom{2p}p \equiv 2 \pmod p$ para todo \omterm{def:b1:arith:prime}{primo} $p$.
  \item Calcule $v_3\binom{1000}{500}$ duas vezes: uma por Kummer
        (escreva $500$ na base $3$ e conte os transportes em $500 +
        500$), outra pela forma digital de Legendre (calcule $s_3(500)$
        e $s_3(1000)$); verifique que as duas dão o mesmo valor.
\end{enumerate}

\medskip
\noindent\textbf{Parte IV --- A paridade do triângulo de Pascal e
uma estimativa de densidade dos \omterm{def:b1:arith:prime}{primos}.}
\begin{enumerate}
  \item Demonstre o critério digital: $\binom nk$ é \emph{ímpar} se,
        e somente se, todo algarismo binário de $k$ é no máximo o
        algarismo correspondente de $n$. Enuncie e demonstre o
        critério análogo para $p \nmid \binom nk$ na base $p$.
  \item Deduza que a linha $n$ do triângulo de Pascal contém exatamente
        $2^{s_2(n)}$ entradas ímpares; verifique nas linhas $4$ e $5$.
  \item Deduza que todas as entradas interiores $\binom nk$ ($0 < k < n$)
        são pares se, e somente se, $n$ é uma potência de $2$.
  \item Demonstre que toda potência de \omterm{def:b1:arith:prime}{primo} que \omterm{def:b1:arith:divides}{divide} $\binom{m+n}m$ é no
        máximo $m + n$: se $p^a \mid \binom{m+n}m$, então $p^a \leq
        m + n$. \emph{(Quantos termos não nulos pode ter a soma da
        questão 11?)}
  \item Deduza que $\binom{2n}n$ \omterm{def:b1:arith:divides}{divide}
        $\operatorname{lcm}(1, 2, \dots, 2n)$ e combine com
        a estimativa inferior $\binom{2n}n \geq \frac{4^n}{2n+1}$
        (que você demonstrará: a entrada central é a maior
        das $2n + 1$ entradas da linha $2n$) para obter
        \[
        \operatorname{lcm}(1, \dots, 2n) \geq \frac{4^n}{2n+1} :
        \]
        os múltiplos comuns dos primeiros inteiros crescem
        \emph{exponencialmente} --- um primeiro vislumbre quantitativo
        da abundância dos \omterm{def:b1:arith:prime}{primos}.
\end{enumerate}

\medskip
\noindent\textbf{Parte V --- Síntese.}
\begin{enumerate}
  \item Encontre o menor $n$ tal que $n!$ termine em pelo menos
        $2026$ zeros. \emph{(Estime $Z(n) \approx n/4$ e depois
        ajuste usando a fórmula exata.)}
  \item Uma última verificação cruzada: mostre que $7$ \emph{não} \omterm{def:b1:arith:divides}{divide}
        $\binom{100}{50}$, primeiro escrevendo $50$ na base $7$ e
        conferindo que a soma $50 + 50$ não tem transportes, e depois
        calculando $v_7(100!)$ e $v_7(50!)$ com a fórmula de
        Legendre.
  \item Onde exatamente o problema usou: (i) a fatoração
        única; (ii) a decomposição pela divisão euclidiana
        $n = p^k n_1 + n_0$; (iii) um argumento de contagem do
        \cref{ch:b1:counting}? Uma frase para cada.
  \item Síntese, num parágrafo curto: a fórmula de Legendre transforma
        uma questão de \omterm{def:b1:arith:divides}{divisibilidade} em aritmética de algarismos, e
        o teorema de Kummer lê a resposta nos transportes de uma única
        soma --- comente essa tradução, as verificações da
        questão 16 e o que a estimativa da questão 21 sugere
        sobre os \omterm{def:b1:arith:prime}{primos} (o enunciado completo, o teorema dos \omterm{def:b1:arith:prime}{números
        primos}, está muito além deste volume; o análogo polinomial
        do instrumental deste capítulo está no
        \cref{ch:b1:poly}).
\end{enumerate}
\end{problem}
